% TOP_PROBLEM: {"problemId": "problem.asser-finite-spectrum-complement-problem", "canonicalTitle": "Finite Spectrum Problem", "releaseRank": 272, "primaryDomain": "logic_foundations_set_theory", "primaryDomainLabel": "Logic, foundations & set theory", "displayStatus": "open_with_solved_subcases", "releaseStatus": "open_with_solved_subcases"}
% !TeX program = lualatex
% Definition and sources only; this file is not an LLM solution attempt.
\documentclass[11pt]{article}
\usepackage[margin=25mm]{geometry}
\usepackage{fontspec}
\setmainfont{Latin Modern Roman}
\usepackage{amsmath,amssymb,mathtools}
\usepackage{unicode-math}
\setmathfont{Latin Modern Math}
\usepackage{xurl,hyperref}
\hypersetup{colorlinks=true,urlcolor=blue}
\setcounter{secnumdepth}{0}
\setlength{\parindent}{0pt}
\setlength{\parskip}{5pt}
\setlength{\emergencystretch}{3em}
\begin{document}
\section*{\texorpdfstring{Finite Spectrum Problem}{Finite Spectrum Problem}}\label{272}
\subsection{Definitions and mathematical statement}
For a first-order sentence $\varphi$ in a finite vocabulary, define its finite spectrum by
\[
 \operatorname{Spec}(\varphi)=\{n\in{\mathbb{Z}}_{>0}:\exists\text{ a structure }\mathcal A,
 \ |\mathcal A|=n,\ \mathcal A\models\varphi\}.
\]
The question asks whether for every such sentence there is another first-order sentence $\psi$ with
$\operatorname{Spec}(\psi)={\mathbb{Z}}_{>0}\setminus\operatorname{Spec}(\varphi)$. The new sentence may use a different finite vocabulary. This is complementation of the set of \emph{possible cardinalities}; simply negating $\varphi$ does not do this, because structures of a fixed size can satisfy $\varphi$ and $\neg\varphi$ simultaneously in different interpretations.
\par\smallskip\textit{Formulation and concept references:} \href{https://arxiv.org/abs/0907.5495}{[S1]}.

\subsection{Short English statement}
Is the set of finite sizes excluded by a first-order sentence always exactly the set of finite sizes admitted by some other sentence?
\subsection{Sources}
\begin{itemize}
\item[S1]A spectrum hierarchy. Submitted scholarly source, accessed 2026-08-31.
\url{https://arxiv.org/abs/0907.5495}.
\textit{Formulation source.}
\end{itemize}
\end{document}
